\documentclass[11pt]{article}
\usepackage[T1]{fontenc}
\usepackage[margin=1in]{geometry}
\usepackage{amsmath,amssymb,amsthm,array,longtable}
\usepackage[hidelinks]{hyperref}
\newtheorem{theorem}{Theorem}
\newtheorem{lemma}[theorem]{Lemma}
\newtheorem{corollary}[theorem]{Corollary}
\newcommand{\Z}{\mathbb Z}
\newcommand{\Gk}{G_k}
\newcommand{\norm}[1]{\lVert #1\rVert_1}
\newcommand{\proofurl}{https://github.com/the-omega-institute/trureturing/blob/50f2a97bb5e704842f2c32c6095f68a0dc62c779/D5/S3/FiniteGroups/}
\title{Bajnok--Beatty Conjectures 4.1 and 1.6\protect\\
(arXiv:2406.04045v1, Section 4, p.~10; Section 1, p.~3):\protect\\
classification and extremal noncyclic order}
\author{Haobo Ma\thanks{ChronoAI Pte Ltd, Singapore.}\and
Wenlin Zhang\thanks{National University of Singapore; ChronoAI Pte Ltd, Singapore.}}
\date{3 November 2026}
\begin{document}
\maketitle
\begin{abstract}
Bajnok--Beatty Conjecture 4.1 in arXiv:2406.04045v1, Section 4,
p.~10, asks which groups
$\Z/2\Z\times\Z/(2k)\Z$ have a two-element signed spanning set of radius
$s$. We prove its four-case classification for every positive $k,s$.
The decisive restriction occurs for $s=2t$ and $k=2t^2-r$, $0\le r<t$:
the deficit $r$ must be even, and divisible by four when $t$ is odd.
A square-covering lattice argument proves this restriction for arbitrary
generating pairs. An explicit lattice construction proves sufficiency.
The lower ranges follow from elementary integer representations.
Counting coefficient residues also proves Conjecture 1.6 of the same
source (Section 1, p.~3): for $s\ge3$, the maximum noncyclic order is
$2s^2$ if $2s+1$ is prime, and $2s^2+2s-(p^2-1)/2$ otherwise,
where $p$ is the smallest prime divisor of $2s+1$.
\end{abstract}
\noindent\textbf{License.} Creative Commons Attribution 4.0 International
(CC BY 4.0).

\section{Statement and attribution}
Write $\Gk=\Z/2\Z\times\Z/(2k)\Z$, where $k\ge1$. A pair of
distinct elements $a,b\in\Gk$ is \emph{$s$-spanning} if every $g\in\Gk$
has a representation $g=xa+yb$ with integers $x,y$ satisfying
$|x|+|y|\le s$. This is equivalent to the undirected Cayley graph with
generating set $\{a,b\}$ having diameter at most $s$.
All posing locators refer to the original version of Bajnok and Beatty
\cite{BB}, arXiv:2406.04045v1: Conjecture 4.1 is in Section 4, p.~10,
and Conjecture 1.6 is in Section 1, p.~3.

\begin{theorem}[Conjecture 4.1]\label{thm:main}
For positive integers $k,s$, the group $\Gk$ has an $s$-spanning pair
if and only if one of the following conditions holds:
\begin{align*}
&s\text{ is odd}, &&2k\le s^2-1;\\
&s\text{ is even}, &&2k\le s^2-s;\\
&s\equiv0\pmod4,\quad k\text{ is even},
&&s^2-s<2k\le s^2;\\
&s\equiv2\pmod4,\quad k\equiv2\pmod4,
&&s^2-s<2k\le s^2.
\end{align*}
\end{theorem}

The bound $2k\le s^2$, the odd-radius classification, and the even-radius
lower range are already Theorems 4.2, 4.3, and 4.4 of \cite{BB}.
The remaining step is the even-radius interval $s^2-s<2k\le s^2$.
We give the full argument, including the lower ranges, to make the
classification self-contained.
The same coefficient-lattice viewpoint gives the extremal order over
all noncyclic finite abelian groups in Corollary~\ref{cor:global}.
Its lower constructions are Bajnok and Beatty's Theorem 1.7;
the residue count below supplies the matching upper bound.

\section{The relation lattice and the order bound}
Let $f:\Z^2\to\Gk$ be $f(x,y)=xa+yb$, for a spanning pair, and let
$K=\ker f$. Then $f$ is surjective. This is the quotient-lattice description
of abelian Cayley diameter in Dougherty and Faber \cite[Proposition 2]{DF}:
the coefficient diamond of radius $s$, translated by $K$, covers $\Z^2$.
Reduction modulo $2\Gk$ gives an
isomorphism $\Z^2/2\Z^2\to\Gk/2\Gk$, since both groups have four
elements. Consequently
\begin{equation}\label{eq:parity}
K\subseteq2\Z^2,\qquad K\not\subseteq4\Z^2.
\end{equation}
For the second assertion, choose $z\in\Z^2$ with $f(z)=(1,0)$.
Then $2z\in K$. If $K\subseteq4\Z^2$, then $z\in2\Z^2$,
contradicting the first coordinate of $f(z)$.

Put $T(u,v)=(u+v,u-v)$ and $M=T^{-1}K$. The image of $T$ consists of
the integer vectors whose coordinates have the same parity. It contains
$K$, and $T$ is injective. Thus
\begin{equation}\label{eq:index}
T(M)=K,\qquad [\Z^2:M]=[\Z^2:K]/2=2k.
\end{equation}
Every $(u,v)\in M$ has $u+v$ even by \eqref{eq:parity}.

\begin{lemma}\label{lem:square}
If $f$ maps the integer diamond $|x|+|y|\le s$ onto $\Gk$, every
coset of $M$ has a representative in $\{0,\ldots,s-1\}^2$.
In particular, $2k\le s^2$.
\end{lemma}
\begin{proof}
First let $s=2t$. Given $z\in\Z^2$, put
$v=T(z-(t,t))+(1,0)$ and choose $w$ with $f(w)=f(v)$ and
$\norm{w}\le2t$. Since $w-v\in K\subseteq2\Z^2$, the integers
$w_1+w_2$ and $w_1-w_2$ are odd. Their absolute values are at most
$2t-1$. Hence
\[
q=\big((w_1+w_2-1)/2,(w_1-w_2-1)/2\big)
\in\{-t,\ldots,t-1\}^2.
\]
As $w=T(q)+(1,0)$, we have $z-(q+(t,t))\in M$, as required.

If $s=2t+1$, instead take $v=T(z-(t,t))$. Now $w_1+w_2$ and
$w_1-w_2$ are even, with absolute values at most $2t$.
Take $q=((w_1+w_2)/2,(w_1-w_2)/2)\in\{-t,\ldots,t\}^2$.
The same subtraction gives the required representative. The square has
$s^2$ points, so \eqref{eq:index} proves the bound.
\end{proof}

\section{Rigidity of a nearly full square}
For a finite-index lattice $L\subseteq\Z^2$, its index is the absolute
determinant of any basis. In particular, the index divides the nonzero
determinant of any two vectors of $L$.

\begin{lemma}\label{lem:rigidity}
Let $n\ge2$, $0\le\delta<n$, and $[\Z^2:L]=n^2-\delta$.
Suppose every coset of $L$ meets $\{0,\ldots,n-1\}^2$.
Then for some $0\le A,B<n$ and $\epsilon\in\{1,-1\}$,
\[
AB=\delta,\qquad
L=\Z(n,-\epsilon B)+\Z(-A,\epsilon n).
\]
\end{lemma}
\begin{proof}
Write $d=n^2-\delta$. Then $d>n(n-1)$ and $2d>n^2$.
For each $0\le y<n$, applying the covering assumption to $(n,y)$
produces a vector $(p,q)\in L$ with
\begin{equation}\label{eq:column}
1\le p\le n,\qquad -(n-1)\le q\le n-1,\qquad 0\le y-q<n.
\end{equation}
First suppose $(p,0)\in L$ for some $1\le p\le n$.
Covering $(0,n)$ produces $(-x,n-z)\in L$ with $0\le x,z<n$.
Their determinant $p(n-z)$ is a positive multiple of $d$ at most $n^2$,
so it equals $d$. If either $p<n$ or $z>0$, this determinant is at most
$n(n-1)$, a contradiction. Thus $p=n$, $z=0$, $\delta=0$, and the
two vectors $(n,0),(-x,n)$ give the desired basis.

Otherwise no vector $(p,0)$ with $1\le p\le n$ belongs to $L$.
Let $b$ be the least positive integer such that $(p,-b)\in L$ for
some $1\le p\le n$, and let $c$ be the least positive integer such
that $(r,c)\in L$ for some $1\le r\le n$. Such integers exist by
\eqref{eq:column} with $y=0,n-1$, and $b,c<n$.
For any $0\le y<n$, the vector in \eqref{eq:column} has either
$q\le-b$, giving $y\le n-1-b$, or $q\ge c$, giving $y\ge c$.
There can be no missing integer between these intervals, so $b+c\le n$.
The determinant of $(p,-b),(r,c)$ is
\[
0<pc+rb\le n(b+c)\le n^2.
\]
It is a multiple of $d$, and therefore equals $d$.
If $b+c<n$, or if both $p<n$ and $r<n$, it is at most $n(n-1)$,
again impossible. Hence $b+c=n$, and either $p=n$ or $r=n$.
If $p=n$, use $(n,-b)$ and $(r,c)-(n,-b)=(-(n-r),n)$.
Here $A=n-r$, $B=b$, and $AB=\delta$.
If $r=n$, use $(n,c)$ and $(p,-b)-(n,c)=(-(n-p),-n)$,
with $A=n-p$, $B=c$, and $\epsilon=-1$.
In each case the two vectors lie in $L$ and their absolute determinant
is $d$, so their span equals $L$.
\end{proof}

\begin{lemma}[Necessary congruences]\label{lem:necessity}
If $s=2t$, $k=2t^2-r$, and $0\le r<t$, any $s$-spanning pair in
$\Gk$ satisfies $2\mid r$; if $t$ is odd, it also satisfies $4\mid r$.
\end{lemma}
\begin{proof}
Apply Lemma~\ref{lem:rigidity} to $M$, using Lemma~\ref{lem:square},
with $n=2t$ and $\delta=2r$. It gives the basis
$(2t,-\epsilon B),(-A,\epsilon2t)$ with $AB=2r$.
Every vector of $M$ has even coordinate sum. Thus $A=2u$ and $B=2v$
for nonnegative integers $u,v$, and $r=2uv$ is even.
If $t$ is odd and $4\nmid r$, both $u,v$ are odd. Applying $T$ to
the two basis vectors gives
\[
2(t-\epsilon v,t+\epsilon v),\qquad
2(-u+\epsilon t,-u-\epsilon t).
\]
Both belong to $4\Z^2$. Since they generate $T(M)=K$, this contradicts
\eqref{eq:parity}.
\end{proof}

\section{A sufficient family}
\begin{lemma}\label{lem:sufficiency}
Let $t\ge1$ and $0\le j<t$. If $t$ is even or $j$ is even, the group
$\Z/2\Z\times\Z/(4(t^2-j))\Z$ has a $2t$-spanning pair.
\end{lemma}
\begin{proof}
Let $L$ be generated by
\[
V=(2(t-j),-2(t+j)),\qquad W=(2(t-1),2(t+1)).
\]
The absolute determinant is $8(t^2-j)$. The greatest common divisor of
the four entries is $2$: after division by $2$ any common divisor
divides both $t-1$ and $t+1$, hence divides $2$; an odd entry occurs
when $t$ is even, or when $t$ is odd and $j$ is even.
The elementary integer row and column reductions for a rank-two matrix
therefore give
\[
\Z^2/L\simeq\Z/2\Z\times\Z/(4(t^2-j))\Z.
\]
Here we use the Smith normal form rule: its first invariant is the
greatest common divisor of the entries and its second is the absolute
determinant divided by the first.

We show that every coset has a representative of $\ell^1$-norm at most
$2t$. Under $T$, the two relation vectors are
\[
T(W)=(4t,-4)=U,\qquad T(V)=(-4j,4t)=V'.
\]
Assume first $t>1$. Their sum has both coordinates positive, since
$j<t$. For a given $z\in\Z^2$, the coset of $T(z)+(2t,2t)$
modulo $\Z U+\Z V'$ therefore has a nonnegative representative.
Choose one, $(X,Y)$, with least coordinate sum.
If $X\ge4t$, subtraction of $U$ preserves nonnegativity and reduces
the sum by $4(t-1)>0$. If $Y\ge4t$, subtraction of $V'$ preserves
nonnegativity and reduces the sum by $4(t-j)>0$.
Consequently $0\le X,Y<4t$.
For $t=1$ necessarily $j=0$; reduction first by $U=(4,-4)$ and then
by $V'=(0,4)$ gives the same square by integer division.

Subtracting $(2t,2t)$ yields a vector $T(w)$ with $w\in z+L$ and
both coordinates in $[-2t,2t-1]$. The identity
\[
|w_1|+|w_2|=\max\{|w_1+w_2|,|w_1-w_2|\}
\]
gives the required norm bound. Finally, $L\subseteq2\Z^2$, so the
images of $(1,0)$ and $(0,1)$ in $\Z^2/L$ are distinct.
Their images under the displayed isomorphism are the desired pair.
\end{proof}

\section{The lower ranges and completion}
For completeness we establish the positive ranges of Theorems 4.3 and
4.4 of \cite{BB} directly.

\begin{lemma}\label{lem:lower}
If $s$ is odd and $2k\le s^2-1$, or if $s$ is even and
$2k\le s^2-s$, the group $\Gk$ has an $s$-spanning pair.
\end{lemma}
\begin{proof}
For $t\ge0$, the two vectors $(1,t),(1,t+1)$ in $\Z/2\Z\times\Z$
represent every $(\eta,z)$ with $|z|\le2t(t+1)$ using coefficient
norm at most $2t+1$. Here is an explicit verification.
For $\eta=0$, write
\[
z=(2t+1)q+w,\quad -t\le w\le t,\quad
q=\left\lfloor\frac{z+t}{2t+1}\right\rfloor.
\]
The bound on $z$ gives $-t\le q\le t$.
Take $x=q-w$, $y=q+w$. Then $x+y$ is even,
$tx+(t+1)y=z$, and $|x|+|y|=2\max(|q|,|w|)\le2t$.
For $\eta=1$, write $z=(2t+1)q+u$, $0\le u\le2t$.
Now $-t-1\le q\le t$. Set
\[
x=q+1-u+t,\qquad y=q+u-t.
\]
Then $x+y=2q+1$ is odd and $tx+(t+1)y=z$.
Also $x-y=2t+1-2u$; thus the same norm identity gives
$|x|+|y|\le2t+1$.
For odd $s=2t+1$, choose an integer representative $z$ of the second
coordinate with $|z|\le k\le2t(t+1)$, and reduce this construction
modulo $2k$. The pair is distinct since its second coordinates differ
by $1$ modulo $2k$.

For even $s=2t$, use $(1,t-1),(1,t)$. The preceding construction
with parameter $t-1$ covers both first coordinates for
$|z|\le B=2t(t-1)$ with coefficient norm at most $2t-1$.
The three integer intervals $ct+[-B,B]$, $c=-1,0,1$, cover
$[-B-t,B+t]$: their endpoints are adjacent or overlap because
$t\le2B+1$. This last inequality holds at $t=1$, and for $t>1$ it
follows from $(t-1)(4t-1)\ge0$.
Choose $c$ with $|z-ct|\le B$, represent $(\eta-c,z-ct)$ in the
smaller range, and add $c(1,t)$. The coefficient norm is at most
$2t$. As $k\le B+t=2t^2-t$, reduction modulo $2k$ gives the claim,
again with distinct elements.
\end{proof}

\begin{proof}[Proof of Theorem~\ref{thm:main}]
Lemma~\ref{lem:square} gives $2k\le s^2$ for every spanning pair.
For odd $s$, the even integer $2k$ is at most $s^2-1$, and
Lemma~\ref{lem:lower} gives the converse.
For even $s=2t$, Lemma~\ref{lem:lower} treats $2k\le s^2-s$.
In the remaining interval put $r=2t^2-k$; then $0\le r<t$.
Lemma~\ref{lem:necessity} says that $r$ is even, and that
$4\mid r$ when $t$ is odd. If $t$ is even this is exactly the
condition that $k$ is even. If $t$ is odd it is exactly
$k\equiv2\pmod4$.
Conversely these congruences give $r=2j$ with $0\le2j<t$ and
$t$ even or $j$ even. Then $k=2(t^2-j)$, and
Lemma~\ref{lem:sufficiency} supplies a $2t$-spanning pair.
\end{proof}

\begin{corollary}\label{cor:maximum}
Among groups $\Gk$ with an $s$-spanning pair, the largest order is
$2s^2$ for even $s\ge2$, and $2s^2-2$ for odd $s\ge3$.
There is no such group for $s=1$.
\end{corollary}
\begin{proof}
The order is $4k$. For odd $s\ge3$, take $2k=s^2-1$.
For even $s$, take $k=s^2/2$: it is even if $s\equiv0\pmod4$,
and it is $2$ modulo $4$ if $s\equiv2\pmod4$.
The upper bounds and the assertion for $s=1$ follow from
Theorem~\ref{thm:main}.
\end{proof}

\begin{corollary}[Even first invariant factor]\label{cor:even-factor}
Let $c\ge2$ be even and $k,s\ge1$. If
$G=\Z/c\Z\times\Z/(ck)\Z$ has an $s$-spanning pair, then
$|G|\le2s^2$. Equality requires $c\mid s$.
In particular, if $s$ is odd then $|G|\le2s^2-2$.
\end{corollary}
\begin{proof}
Let $f:\Z^2\to G$ be the surjection induced by the pair, and put
$K=\ker f$. The induced surjection $(\Z/c\Z)^2\to G/cG$ is bijective,
since both groups have order $c^2$. Thus $K\subseteq c\Z^2\subseteq2\Z^2$.
For $M=T^{-1}K$ we have $T(M)=K$ and $[\Z^2:M]=|G|/2$.
The proof of Lemma~\ref{lem:square}, which uses only $K\subseteq2\Z^2$,
therefore gives $|G|/2\le s^2$.
If equality holds, $s\ge2$ and Lemma~\ref{lem:rigidity} applies with
$n=s$, $\delta=0$. Its basis has $AB=0$. If $B=0$, applying $T$ to
the first basis vector gives $(s,s)\in K$. If $A=0$, applying $T$ to
the second gives $(\epsilon s,-\epsilon s)\in K$.
In either case $K\subseteq c\Z^2$ implies $c\mid s$.
For odd $s$ equality is impossible; since $4\mid|G|$,
the inequality sharpens to $|G|\le2s^2-2$.
\end{proof}

This extends the even-$c$ part of Theorems 1.8 and 3.2(1) of \cite{BB}:
the source assumes an $s$-spanning pair of the special form
$\{(1,u),(1,v)\}$, whereas Corollary~\ref{cor:even-factor} permits any pair.

\section{Residue counts and Conjecture 1.6}
\begin{lemma}[Odd first invariant factor]\label{lem:odd-factor}
Let $c\ge3$ be odd, and write $s=ct+r$ with $0\le r<c$.
If $G=\Z/c\Z\times\Z/(ck)\Z$ has an $s$-spanning pair, then
\[
k\le F_c(s):=
\begin{cases}
2t^2,&0\le r<(c-1)/2,\\
2t(t+1),&(c-1)/2\le r<c-1,\\
2t(t+1)+1,&r=c-1.
\end{cases}
\]
No regularity hypothesis on the pair is required.
\end{lemma}
\begin{proof}
For the surjection $f:\Z^2\to G$ induced by the pair, put $K=\ker f$.
As in Corollary~\ref{cor:even-factor}, $K\subseteq c\Z^2$.
Consequently all representatives of any one element of $G$ have the
same coefficient residue modulo $c$. Each of the $c^2$ coefficient
residues corresponds to one coset of $cG$, with $k$ elements.
Each residue class in the diamond $|x|+|y|\le s$ must therefore
contain at least $k$ points.

Put $U=x+y$ and $V=x-y$. Since $c$ is odd, this is an invertible
change of coordinates modulo $c$. Over the integers it identifies
the diamond with $|U|,|V|\le s$, $U\equiv V\pmod2$.
For a residue $i$ modulo $c$, let $E_i,O_i$ count its even and odd
representatives in $[-s,s]$. A prescribed pair of residues for $U,V$
therefore contains $E_iE_j+O_iO_j$ diamond points.

Choose $i=s+1\pmod c$. Its representatives are $s+1-\ell c$,
$1\le\ell\le\lfloor(2s+1)/c\rfloor$, with alternating parity.
If $r<(c-1)/2$, there are $2t$ representatives, so $E_i=O_i=t$.
Taking $j=i$ gives $2t^2$ points.
If $(c-1)/2\le r<c-1$, both $i=s+1$ and $j=s+2$ have $2t+1$
representatives: for the latter the number is
$\lfloor(2s+2)/c\rfloor=2t+1$.
Their first representatives have opposite parity. Thus their counts
are $(t,t+1)$ and $(t+1,t)$, giving $2t(t+1)$ points.
If $r=c-1$, take $j=i$ again. There are $2t+1$ representatives, giving
$t^2+(t+1)^2=2t(t+1)+1$ points. In each case $k$ is at most this count.
\end{proof}

\begin{corollary}[Conjecture 1.6]\label{cor:global}
For $s\ge3$, the largest order $N(s)$ of a noncyclic finite abelian
group with an $s$-spanning pair is
\[
N(s)=\begin{cases}
2s^2,&2s+1\text{ is prime},\\
2s^2+2s-(p^2-1)/2,&2s+1\text{ is composite},
\end{cases}
\]
where $p$ is the smallest prime divisor of $2s+1$ in the second case.
\end{corollary}
\begin{proof}
Such a group has the form $G=\Z/c\Z\times\Z/(ck)\Z$ with $c\ge2$.
Even $c$ gives $|G|\le2s^2$ by Corollary~\ref{cor:even-factor}.
For odd $c$, apply Lemma~\ref{lem:odd-factor}.
In its first case, $c^2F_c(s)=2c^2t^2\le2s^2$.
In its middle case,
\[
c^2F_c(s)-2s^2=2ct(c-2r)-2r^2.
\]
This is negative for $r>(c-1)/2$. At $r=(c-1)/2$ it gives exactly
\[
|G|\le2s^2+2s-(c^2-1)/2,\qquad c\mid 2s+1.
\]
In the last case the difference is
\[
c^2F_c(s)-2s^2=-2ct(c-2)-c^2+4c-2<0.
\]
For $c\ge5$ the constant term is already negative; for $c=3$, $s\ge3$
and $r=2$ imply $t\ge1$, giving $1-6t<0$.

If $2s+1$ is prime, its only possible divisor $c\ge3$ in the exceptional
middle case is $c=2s+1$. Then $t=0$ and $F_c(s)=0$, incompatible with
$k\ge1$. Hence $|G|\le2s^2$.
If $2s+1$ is composite, every exceptional $c$ satisfies $c\ge p$;
its bound is at most $2s^2+2s-(p^2-1)/2$.
This exceeds $2s^2$: writing $2s+1=pq$ with $q\ge p$ gives excess
$pq-1-(p^2-1)/2\ge(p^2-1)/2>0$.

The matching constructions are those of \cite[Theorem 1.7,
Proposition 2.3 and Corollary 2.5]{BB}:
$\Z/s\Z\times\Z/(2s)\Z$ attains $2s^2$.
In the composite case, put $q=(2s+1)/p$ and $t=(q-1)/2\ge1$.
The pair $\{(1,t),(1,t+1)\}$ in
$\Z/p\Z\times\Z/(2pt(t+1))\Z$ is $s$-spanning by
\cite[Theorem 2.4]{BB}, with order $2s^2+2s-(p^2-1)/2$.
All these groups are noncyclic and the indicated pairs are distinct.
\end{proof}

\begin{corollary}[Theorem 3.3 without regularity]\label{cor:universal}
For $s\ge3$, every noncyclic finite abelian group with an $s$-spanning
pair has order at most $2s^2+2s-4$.
Equality holds precisely for the isomorphism type
$\Z/3\Z\times\Z/((2s^2+2s-4)/3)\Z$, with $s\equiv1\pmod3$.
\end{corollary}
\begin{proof}
The nonexceptional bounds above are at most $2s^2<2s^2+2s-4$.
In the exceptional bound equality requires $c=3$, hence
$s\equiv1\pmod3$; the construction in Corollary~\ref{cor:global}
attains it. This also extends Theorem 1.8 to arbitrary pairs:
the odd-$c$ residue bounds imply all its stated upper bounds, while
Corollary~\ref{cor:even-factor} treats even $c$.
\end{proof}

\paragraph{Small radii.}
There is no noncyclic group with a $1$-spanning pair:
even $c$ would give $c^2\le|G|\le2$, and odd $c\ge3$ gives
$k\le F_c(1)=0$. At radius $2$, the even-$c$ bound gives order at
most $8$, while the odd-$c$ bound permits only $c=3$, $k=1$;
$\{(1,0),(0,1)\}$ spans $\Z_3^2$ with radius $2$.
Thus $N(2)=9$. In particular the first bound in
\cite[Theorem 1.7]{BB} requires $s\ge2$, and its Theorem 1.9
requires the radius-one exception stated in Corollary~\ref{cor:maximum}.

\paragraph{Perfect spanning sets (Conjecture 1.10, Section 1, p.~4).}
For $m=1$, perfection forces a cyclic group of order $2s+1$ with its
chosen generator. For $s=1$, perfection is exactly the disjoint union
$G=\{0\}\sqcup A\sqcup(-A)$. These are the boundary cases already
stated in \cite{BB}. For $m=2$, a perfect set has
$|G|=2s^2+2s+1$, hence $G$ is cyclic by Corollary~\ref{cor:universal}
when $s\ge3$. When $s=2$, Lemma~\ref{lem:odd-factor} and the even-$c$
bound give maximum noncyclic order $9$, attained by $\Z_3^2$;
when $s=1$, the order is the prime $5$. Thus cyclicity holds for all $s$.

The printed multiplier domain in Conjecture 1.10 needs correction:
$d\in\Z_{2s+1}$ should be $d\in(\Z_{2s^2+2s+1})^\times$.
Indeed, for $s=2$, $A=\{2,10\}\subset\Z/13\Z$ is perfect:
$0,\pm2,\pm10,\pm4,\pm20,\pm12,\pm8$ give all thirteen residues.
It equals $5\{2,3\}$, but equals none of $d\{2,3\}$ for
$0\le d\le4$. This corrects the domain; the full two-generator
shape assertion with that corrected domain is not proved here.

\section{Coverage and verification}
The following table separates the exact posing statement from other items
in the same source. The term ``not covered'' refers to the scope of this note.
\begin{center}\small
\begin{tabular}{p{.26\linewidth}p{.65\linewidth}}
\hline
Source item in \cite{BB}&Coverage\\ \hline
Conjecture 4.1 (p.~10)&Proved: Theorem~\ref{thm:main}, all four clauses and every positive $k,s$.\\
Theorems 4.2--4.4&Prior proofs cited; Lemmas~\ref{lem:square} and \ref{lem:lower}
give complete proofs of the order bound and lower ranges.\\
Theorems 1.7, 1.9&Prior proofs cited for $s\ge2$;
refuted-and-corrected radius-one maximum/lower-bound assertions:
no noncyclic group has a $1$-spanning pair.\\
Theorems 1.8, 3.2&Proved without regularity: Corollary~\ref{cor:even-factor}
and Lemma~\ref{lem:odd-factor} imply their order bounds for arbitrary pairs.\\
Conjecture 1.6 (p.~3)&Proved: Corollary~\ref{cor:global}, every $s\ge3$;
the lower constructions are prior proofs cited from Theorem 1.7.\\
Theorem 3.3&Proved without regularity: Corollary~\ref{cor:universal},
including its equality isomorphism type.\\
Conjecture 1.10 (p.~4)&Refuted-and-corrected multiplier domain;
proved cyclicity for $m=2$ and the elementary $m=1$ and $s=1$ cases.
Corrected $m=2$ shape not covered. For $m\ge3,s\ge2$, prior theorem
statement cited from \cite[Theorem 1.1]{CHL}.\\
Sections 1--3, other items&Prior proofs cited: Proposition 1.1,
Theorems 1.2--1.5, Lemma 2.1, Propositions 2.2--2.3,
Theorem 2.4, Corollary 2.5 and Theorem 3.1.
Directed questions for $m\ge3$ and exact orders with $m\ge3$
generators are not covered.\\
\cite{BBJ}, Questions 6, 7, 10--12&Not covered: restricted unsigned bases
have a different coefficient set from the signed two-generator diamonds.\\ \hline
\end{tabular}
\end{center}

For the Conjecture 1.10 comparison, a perfect $s$-spanning set with $m$
elements makes the coefficient map $\Z^m\to G$ bijective on the integer
Lee ball $B_m(s)=\{x:\sum_i|x_i|\le s\}$. Its kernel therefore gives a
lattice tiling by that ball. The exclusion for $m\ge3$, $s\ge2$ is the
statement of Chen, Huang, and Luo \cite[Theorem 1.1]{CHL}, rather than
a consequence of our two-dimensional argument. Their lattice-tiling
preprint was posted on 25 October 2026; their subsequent same-day
preprint \cite[Theorem 1.1]{CHLA} states the stronger exclusion for
arbitrary centre sets. These statements concern existence parameters,
not the individual two-dimensional perfect sets.

The pinned Lean declarations have the following scopes:
\begin{itemize}
\item \href{\proofurl BajnokBeattyAllRegimes.lean}{\texttt{BajnokBeattyAllRegimes.result}}:
Theorem~\ref{thm:main}, with positive natural $k,s$, distinct $a,b$,
and the signed coefficient bound.
\item \href{\proofurl BajnokBeattyCayleyNecessity.lean}{\texttt{BajnokBeattyCayleyNecessity.}\allowbreak\texttt{top\_strip\_necessity}}:\newline
Lemma~\ref{lem:necessity} for arbitrary pairs, without requiring distinctness.
\item \href{\proofurl BajnokBeattyCayleySufficiency.lean}{\texttt{BajnokBeattyCayleySufficiency.}\allowbreak\texttt{result}}:
Lemma~\ref{lem:sufficiency}, including $t=1,j=0$.
\item \href{\proofurl NearSquareLatticeRigidity.lean}{\texttt{NearSquareLatticeRigidity.}\allowbreak\texttt{near\_square\_rigidity}}:
the lattice statement of Lemma~\ref{lem:rigidity}.
\end{itemize}
All four links select commit
\texttt{50f2a97bb5e704842f2c32c6095f68a0dc62c779}.
The arguments above are written proofs of the stated lemmas and corollaries.
Corollary~\ref{cor:even-factor} is a written consequence of the square and
rigidity lemmas. Lemma~\ref{lem:odd-factor},
Corollaries~\ref{cor:global}--\ref{cor:universal}, and the assertions about
Conjecture 1.10 have written proofs or the displayed literature attribution;
the four listed declarations do not separately state these results.

The accompanying \texttt{verify.py} independently computes minimum diameters
over all distinct two-element subsets of $\Gk$ for $1\le k\le40$,
and compares the classification for $1\le s\le12$.
It also checks the displayed sufficient family through $t=12$ and the
integer formulas for the lower ranges through $t=12$.
It also exhausts all distinct pairs for $c\in\{4,6\}$, $1\le k\le4$
in the even-first-invariant-factor corollary.
It compares the odd-$c$ residue bounds with explicit diamond counts,
checks the extremal formula over bounded radii, and verifies the displayed
perfect-set counterexample and the attaining pairs.
An exhaustive comparison of all noncyclic two-generator abelian groups
of order at most $61$ gives the global maxima for $2\le s\le5$.
These finite checks supplement, rather than replace, the proofs for
unbounded parameters.

\begin{thebibliography}{5}
\bibitem{BB}
B. Bajnok and W. K. Beatty,
\emph{On the Diameter of Undirected Cayley Graphs of Finite Abelian Groups},
arXiv:2406.04045v1, 6 June 2024,
\href{https://arxiv.org/abs/2406.04045v1}{arXiv:2406.04045v1}.
Published as \emph{Combinatorics and Number Theory} \textbf{14} (2025),
no.~1, 1--12,
\href{https://doi.org/10.2140/cnt.2025.14.1}{doi:10.2140/cnt.2025.14.1}.
\bibitem{DF}
R. Dougherty and V. Faber,
\emph{The degree-diameter problem for several varieties of Cayley graphs,
I: The Abelian case},
\emph{SIAM Journal on Discrete Mathematics} \textbf{17} (2004), no.~3,
478--519,
\href{https://doi.org/10.1137/S0895480100372899}{doi:10.1137/S0895480100372899}.
The cited Proposition 2 is in
\href{https://arxiv.org/abs/math/9412223v2}{arXiv:math/9412223v2},
8 September 2000, p.~6.
\bibitem{CHL}
B. Chen, J. Huang, and G. Luo,
\emph{The Golomb--Welch conjecture for lattice tilings},
\href{https://arxiv.org/abs/2610.33983v1}{arXiv:2610.33983v1},
25 October 2026, Theorem~1.1.
\bibitem{CHLA}
B. Chen, J. Huang, and G. Luo,
\emph{The Golomb--Welch conjecture for arbitrary centre sets},
\href{https://arxiv.org/abs/2610.34291v1}{arXiv:2610.34291v1},
25 October 2026, Theorem~1.1.
\bibitem{BBJ}
B. Bajnok, C. Berson, and H. A. Just,
\emph{On Perfect Bases in Finite Abelian Groups},
\href{https://arxiv.org/abs/2211.13675v1}{arXiv:2211.13675v1},
24 November 2022, Section~4.
\end{thebibliography}
\end{document}
